\chapter{Mehrere Zufallsvariablen}

Bevor wir uns anschauen was passiert, wenn wir mehrere Zufallsvariablen haben, schauen wir uns 
zuerst bedingte Zufallsvariablen an.

\begin{definition}
    Sei $X$ eine Zufallsvariable auf dem Wahrscheinlichkeitsraum $\Omega$ mit $A \subseteq \Omega$
    und Pr$[A] > 0$. Die \textbf{bedingte Zufallsvariable} $X | A$ ist dieselbe Funktion wie $X$,
    aber der Definitionsbereich ist auf die Teilmenge $A$ eingeschränkt. \\

    \textbf{Dichtefunktion}

    $$f_{X | A}: \mathbb{R} \rightarrow [0, 1], x \mapsto \text{Pr}[X = x | A]$$

    \textbf{Verteilungsfunktion}

    $$F_{X | A}: \mathbb{R} \rightarrow [0, 1], x \mapsto \text{Pr}[X \leq x | A]$$

    \textbf{Erwartungswert}

    $$\text{E}[X | A] := \sum_{x \in W_X} x \cdot \text{Pr}[X = x | A] = \frac{1}{\text{Pr}[A]} \cdot \sum_{\omega \in A} X(\omega) \cdot \text{Pr}[\omega]$$

\end{definition}
\bigskip

Oftmals interessieren wir uns nun für mehrere Zufallsvariablen zugleich. In diesem Kapitel beschäftigen 
wir uns mit der Frage, wie wir mit mehreren Zufallsvariablen über demselbem Wahrscheinlichkeitsraum rechnen 
können - selbst wenn diese sich gegenseitig beeinflussen. Dazu definieren wir:

\begin{definition}
    Für zwei Zufallsvariablen $X, Y$ heisst die Funktion
    $$f_{X,Y}(x, y) := \text{Pr}[X = x, Y = y]$$
    die \textbf{gemeinsame Dichte} der Zufallsvariable $X$ und $Y$. Von der gemeinsamen Dichte kann man
    auch wieder zu der Dichte der jeweiligen Zufallsvariablen übergehen - den sogenannten Randdichten:
    $$f_X(x) = \sum_{y \in W_y} f_{X,Y}(x,y) \text{   und   } f_Y(y) = \sum_{x \in W_X} f_{X,Y}(x,y)$$
\end{definition}
\bigskip

Für ein Beispiel siehe Skript. Analog dazu kann man auch gemeinsame Verteilung und Randverteilung definieren, 
was in der Vorlesung jedoch nicht verwendet wurde.

\section{Unabhängigkeit von Zufallsvariablen}

Analgo zu Ereignissen können wir auch für Zufallsvariablen Unabhängigkeit definieren. Dies ist offensichtlich
der Fall, wenn Zufallsvariablen physikalisch voneinander unabhängig sind. Es kann aber auch nur implizit
durch nachrechnen bestätigt werden.

\begin{definition}
    Zufallsvariablen $X_1, ... , X_n$ heissen unabhängig genau dann, wenn für alle $(x_1, ... , x_n) \in W_{X_1} \times ... \times W_{X_n}$ 
    gilt
    $$\text{Pr}[X_1 = x_1, ... , X_n = x_n] = \text{Pr}[X_1 = x_1] \cdot ... \cdot \text{Pr}[X_n = x_n]$$
    Oder alternativ
    $$f_{x_1, ... , x_n}(x_1, ... , x_n) = f_{X_1}(x_1) \cdot ... \cdot f_{X_n}(x_n)$$
\end{definition}
\bigskip

Aus dieser Definition können wir direkt folgern, dass zwei Indikatorvariablen unabhängig sind, wenn ihre
Ereignissen unabhängig sind. Allgemein gilt für zwei Indikatorvariablen:
$$X \text{ und } Y \text{ sind unabhängig } \Leftrightarrow f_{X,Y}(1,1) = f_X(1) \cdot f_Y(1)$$
Aus der Definition folgen einige Lemmas:
